\begin{abstract}

% The emergence of CXL 3.0's fabric-attached memory enables multiple hosts to share memory resources transparently, creating new opportunities for distributed computing workloads in machine learning and data analytics. However, efficient synchronization remains a critical challenge for processes collaborating through shared CXL memory. This paper presents the first comprehensive evaluation of mutex implementations across CXL's two distinct coherence models: software-managed coherence using load/store operations with fences, and hardware-managed coherence leveraging read-modify-write atomics with back-invalidation snooping. We implement and evaluate a diverse set of synchronization primitives on a BittWare IA-780i FPGA platform with Intel CXL IP, introducing the first FPGA-based implementation of CXL's hardware coherence protocol with back-invalidation support. Our evaluation framework includes both classical software locks (Bakery, Lamport's Fast Algorithm, Peterson, Knuth, and novel Elevator Locks ) that rely solely on memory ordering guarantees, and hardware locks (spinlocks, ticket, MCS) that exploit atomic operations. Through configurable delay buffers, we characterize performance across latencies ranging from 400ns to 10 us, simulating deployment scenarios from local memory expansion to cross-rack fabric access. Our results reveal critical trade-offs between software and hardware synchronization approaches in CXL environments. Software locks demonstrate predictable scaling characteristics independent of cache coherence traffic but suffer from O(N log N) memory overhead and higher base latency. Hardware locks supported by hardware coherence achieve superior performance with a small number of threads but degrade with increasing cache line contention. We identify specific crossover points where software-managed coherence becomes competitive, suggesting that applications can strategically avoid the complexity of full hardware coherence for certain synchronization patterns. These findings provide essential guidance for system designers choosing between CXL coherence models and for developers optimizing synchronization primitives for disaggregated memory architectures.

Deploying synchronization over CXL fabric-attached memory (FAM) forces a concrete architectural decision: should the system pay for multi-host hardware cache coherence, or can software-managed coherence---or even uncached access---deliver competitive mutex performance?
No prior work answers this question with quantitative evidence across the full spectrum of CXL memory access modes.
We present the first systematic study that evaluates mutex performance across \emph{four} access modes on real CXL hardware: (1)~cached with hardware coherence (BISnp), (2)~cached with software-managed coherence, (3)~uncached (UC) access, and (4)~non-temporal (NT) streaming access.
Using a BittWare IA-780i FPGA platform with programmable delay injection (400\,ns--10\,$\mu$s), we benchmark 13 lock implementations---spanning classical software locks (Bakery, Peterson, Elevator) and hardware locks (spin, ticket, MCS)---under each access mode and latency configuration.
Our results reveal sharp, quantifiable trade-offs.
Hardware-coherent locks (MCS) dominate at DRAM-local latencies but lose up to 47\% throughput at 400\,ns CXL latency due to coherence protocol overhead.
Uncached software locks close the gap: at 10\,$\mu$s latency, the Elevator-BL lock under UC access achieves within 12\% of MCS under full hardware coherence, eliminating the need for costly coherence infrastructure.
Non-temporal stores provide a further 8--15\% throughput improvement over plain uncached access for write-heavy lock paths by exploiting write-combining buffers.
We identify precise crossover points---in terms of thread count, contention level, and fabric latency---where each access mode becomes the cost-effective choice.
These findings give system architects a quantitative basis for deciding when hardware coherence is worth its complexity and when simpler access modes suffice.

\end{abstract}
